% ============================ WORKSHEET 4 ==================================
\worksheetheading{4}{Electronic configuration and periodic table}
\data{Atomic numbers: H 1, He 2, Li 3, Be 4, B 5, C 6, N 7, O 8, F 9, Ne 10, Na 11, Mg 12, Al 13,
Si 14, P 15, S 16, Cl 17, Ar 18, K 19, Ca 20, Sc 21, Ti 22, V 23, Cr 24, Mn 25, Fe 26, Co 27, Ni 28,
Cu 29, Zn 30, Ga 31, Ge 32, As 33, Se 34, Br 35, Kr 36, Rb 37, Sr 38. No calculator needed.}

\wq{1}{Capacity of subshells}
Complete: an s subshell has \dots\ orbital(s) and holds at most \dots\ electrons; p: \dots\ and
\dots; d: \dots\ and \dots; f: \dots\ and \dots. Which principle limits an orbital to two electrons?
\begin{solution}
s: 1 and 2;\ p: 3 and 6;\ d: 5 and 10;\ f: 7 and 14. The Pauli exclusion principle: two electrons of an
orbital share $n,\ell,m_\ell$, so they must differ by $m_s$ ($+\frac12$ and $-\frac12$).
\end{solution}

\wq{1}{The three rules}
State the Aufbau principle, the Pauli exclusion principle and Hund's rule in one sentence each.
\begin{solution}
Aufbau: in the ground state, electrons fill the lowest-energy orbitals first (Madelung order).
Pauli: two electrons of an atom cannot have the same four quantum numbers (at most two per orbital,
with opposite spins). Hund: in a subshell, electrons occupy separate orbitals with parallel spins
before pairing.
\end{solution}

\wq{1}{Writing configurations}
Write the full ground-state configuration of O, Na, Cl and Ar, then the shorthand form with a noble
gas core.
\begin{solution}
O: $1\mathrm s^2 2\mathrm s^2 2\mathrm p^4$ $= [\mathrm{He}]\,2\mathrm s^2 2\mathrm p^4$.\quad
Na: $1\mathrm s^2 2\mathrm s^2 2\mathrm p^6 3\mathrm s^1 = [\mathrm{Ne}]\,3\mathrm s^1$.\quad
Cl: $1\mathrm s^2 2\mathrm s^2 2\mathrm p^6 3\mathrm s^2 3\mathrm p^5 = [\mathrm{Ne}]\,3\mathrm s^2 3\mathrm p^5$.\quad
Ar: $1\mathrm s^2 2\mathrm s^2 2\mathrm p^6 3\mathrm s^2 3\mathrm p^6 = [\mathrm{Ne}]\,3\mathrm s^2 3\mathrm p^6$.
\end{solution}

\wq{1}{Core and valence}
Give the number of core electrons and of valence electrons of C, Mg, S and Ar.
\begin{solution}
C $[\mathrm{He}]\,2\mathrm s^2 2\mathrm p^2$: 2 core, 4 valence.\quad
Mg $[\mathrm{Ne}]\,3\mathrm s^2$: 10 core, 2 valence.\quad
S $[\mathrm{Ne}]\,3\mathrm s^2 3\mathrm p^4$: 10 core, 6 valence.\quad
Ar $[\mathrm{Ne}]\,3\mathrm s^2 3\mathrm p^6$: 10 core, 8 valence.
\end{solution}

\wq{2}{Box diagrams and Hund's rule}
(a) Draw the box diagram of the valence shell of N, O and F and give the number of unpaired
electrons. (b) Which of these 2p$^3$ diagrams respects Hund's rule? Explain what is wrong with the
others: (i) \obox{d,u,e}\quad (ii) \obox{u,u,u}\quad (iii) \obox{u,w,u}.
\begin{solution}
(a) N: \obox{d}\;\obox{u,u,u}, 3 unpaired.\quad O: \obox{d}\;\obox{d,u,u}, 2 unpaired.\quad
F: \obox{d}\;\obox{d,d,u}, 1 unpaired.\\
(b) (ii). In (i) two electrons pair up while an orbital is still empty; in (iii) the spins are not
all parallel. Both are allowed by Pauli but are excited states.
\end{solution}

\wq{2}{From configuration to element}
Identify the element and give its period, group and block:
(a) $1\mathrm s^2 2\mathrm s^2 2\mathrm p^6 3\mathrm s^2 3\mathrm p^4$;\ (b) $[\mathrm{Ne}]\,3\mathrm s^1$;\
(c) $[\mathrm{Ar}]\,3\mathrm d^6 4\mathrm s^2$;\ (d) $[\mathrm{Kr}]\,5\mathrm s^2$.
\begin{solution}
(a) 16 electrons: S, period 3, group 16, p block.\quad (b) 11: Na, period 3, group 1, s block.\quad
(c) $18 + 8 = 26$: Fe, period 4, group 8, d block.\quad (d) $36 + 2 = 38$: Sr, period 5, group 2, s block.
\end{solution}

\wq{2}{From position to configuration}
Write the shorthand configuration of the element located at: (a) period 3, group 15;
(b) period 4, group 2; (c) period 4, group 17; (d) period 5, group 1. Name each element.
\begin{solution}
(a) $[\mathrm{Ne}]\,3\mathrm s^2 3\mathrm p^3$: phosphorus.\quad (b) $[\mathrm{Ar}]\,4\mathrm s^2$: calcium.\quad
(c) $[\mathrm{Ar}]\,3\mathrm d^{10} 4\mathrm s^2 4\mathrm p^5$: bromine (the 3d subshell of period 4 is full
after group 12).\quad (d) $[\mathrm{Kr}]\,5\mathrm s^1$: rubidium.
\end{solution}

\wq{2}{Ground, excited or impossible?}
Classify each configuration as ground state, excited state or impossible, and name the atom:
(a) $1\mathrm s^2 2\mathrm s^2 2\mathrm p^6 3\mathrm s^1$;\ (b) $1\mathrm s^2 2\mathrm s^2 2\mathrm p^5 3\mathrm s^1$;\
(c) $1\mathrm s^2 2\mathrm s^3$;\ (d) $1\mathrm s^2 2\mathrm p^1$;\ (e) $1\mathrm s^2 2\mathrm s^2 2\mathrm p^7$;\
(f) $1\mathrm s^2 2\mathrm s^2 2\mathrm p^6 3\mathrm d^1$.
\begin{solution}
(a) Ground state of Na (11 e$^-$).\quad (b) Excited Ne (10 e$^-$, one electron promoted from 2p to 3s).\quad
(c) Impossible: an s subshell holds at most 2 electrons (Pauli).\quad (d) Excited Li (3 e$^-$, 2p filled
before 2s).\quad (e) Impossible: p holds at most 6.\quad (f) Excited Na (3d instead of 3s).
\end{solution}

\wq{3}{Ions and noble gases}
Write the configuration of $\ce{Na+}$, $\ce{Mg^2+}$, $\ce{O^2-}$, $\ce{Cl-}$ and $\ce{S^2-}$. Which noble gas
does each resemble? Explain why these ions are the ones usually formed.
\begin{solution}
$\ce{Na+}$ (10 e$^-$), $\ce{Mg^2+}$ (10), $\ce{O^2-}$ (10): $1\mathrm s^2 2\mathrm s^2 2\mathrm p^6 = [\mathrm{Ne}]$.\quad
$\ce{Cl-}$ (18), $\ce{S^2-}$ (18): $[\mathrm{Ne}]\,3\mathrm s^2 3\mathrm p^6 = [\mathrm{Ar}]$.
By losing or gaining a few valence electrons, each atom reaches the filled, very stable configuration of
the nearest noble gas.
\end{solution}

\wq{3}{Iron and its ions}
(a) Write the configuration of Fe ($Z = 26$) and draw the box diagram of its 3d and 4s subshells.
How many unpaired electrons does it have? (b) Write the configuration of $\ce{Fe^2+}$ and $\ce{Fe^3+}$.
(c) $\ce{Fe^3+}$ is particularly stable. Suggest why.
\begin{solution}
(a) $[\mathrm{Ar}]\,3\mathrm d^6 4\mathrm s^2$; 3d: \obox{d,u,u,u,u}, 4s: \obox{d}: 4 unpaired electrons.\\
(b) Remove 4s first: $\ce{Fe^2+}$: $[\mathrm{Ar}]\,3\mathrm d^6$;\ $\ce{Fe^3+}$: $[\mathrm{Ar}]\,3\mathrm d^5$.\\
(c) $3\mathrm d^5$ is a half-filled subshell: five parallel electrons, one per orbital, a particularly
stable arrangement (the same reason as for chromium's exception).
\end{solution}

\wq{3}{Exceptions: chromium and copper}
(a) Write the configuration predicted by Madelung's rule for Cr ($Z = 24$) and Cu ($Z = 29$).
(b) The real configurations are $[\mathrm{Ar}]\,3\mathrm d^5 4\mathrm s^1$ and
$[\mathrm{Ar}]\,3\mathrm d^{10} 4\mathrm s^1$. Explain. (c) How many unpaired electrons do Cr and Cu have?
\begin{solution}
(a) Cr: $[\mathrm{Ar}]\,3\mathrm d^4 4\mathrm s^2$;\ Cu: $[\mathrm{Ar}]\,3\mathrm d^9 4\mathrm s^2$.\\
(b) 3d and 4s are very close in energy. Moving one electron from 4s to 3d gives a half-filled
($\mathrm d^5$) or completely filled ($\mathrm d^{10}$) subshell, which is extra stable; the total energy is
lower.\\
(c) Cr: 3d \obox{u,u,u,u,u}, 4s \obox{u}: 6 unpaired.\quad Cu: 3d full, 4s \obox{u}: 1 unpaired.
\end{solution}

\wq{3}{Four quantum numbers}
(a) Give a possible set $(n,\ell,m_\ell,m_s)$ for each of the five valence electrons of nitrogen in
its ground state. (b) Which of these sets could describe an electron of a ground-state phosphorus atom?
(i) $(3,1,0,+\frac12)$; (ii) $(3,2,0,+\frac12)$; (iii) $(2,1,-1,-\frac12)$; (iv) $(3,0,0,0)$.
\begin{solution}
(a) $(2,0,0,+\frac12)$ and $(2,0,0,-\frac12)$ for 2s$^2$; then, by Hund, $(2,1,-1,+\frac12)$,
$(2,1,0,+\frac12)$, $(2,1,+1,+\frac12)$ for 2p$^3$ (all $-\frac12$ is equally correct).\\
(b) P $= 1\mathrm s^2 2\mathrm s^2 2\mathrm p^6 3\mathrm s^2 3\mathrm p^3$. (i) yes, a 3p electron;
(ii) no, 3d is empty in the ground state; (iii) yes, a core 2p electron; (iv) impossible,
$m_s$ must be $\pm\frac12$.
\end{solution}

\wq{4}{Why 2, 8, 8, 18, 18, 32?}
Using Madelung's order and the fact that each period starts when a new $n$s subshell starts filling,
list the subshells filled in each of the first six periods and deduce the number of elements in each.
\begin{solution}
Period 1: 1s (2).\quad Period 2: 2s 2p ($2+6 = 8$).\quad Period 3: 3s 3p (8; 3d comes after 4s).\quad
Period 4: 4s 3d 4p ($2+10+6 = 18$).\quad Period 5: 5s 4d 5p (18).\quad
Period 6: 6s 4f 5d 6p ($2+14+10+6 = 32$). Period 7 is built the same way (7s 5f 6d 7p: 32).
\end{solution}

\wq{4}{Magnetic atoms}
Atoms with unpaired electrons are attracted by a magnet (paramagnetic); atoms with none are not
(diamagnetic). Find the number of unpaired electrons of Mn ($Z = 25$), Fe, Zn ($Z = 30$), Cu and
$\ce{Zn^2+}$. Which are diamagnetic? Which ion of manganese, $\ce{Mn^2+}$ or $\ce{Mn^3+}$, has the most
unpaired electrons?
\begin{solution}
Mn $[\mathrm{Ar}]\,3\mathrm d^5 4\mathrm s^2$: 5.\quad Fe $3\mathrm d^6 4\mathrm s^2$: 4.\quad
Zn $3\mathrm d^{10} 4\mathrm s^2$: 0.\quad Cu $3\mathrm d^{10} 4\mathrm s^1$: 1.\quad
$\ce{Zn^2+}$ $[\mathrm{Ar}]\,3\mathrm d^{10}$: 0.\\
Diamagnetic: Zn and $\ce{Zn^2+}$.\quad $\ce{Mn^2+}$ $[\mathrm{Ar}]\,3\mathrm d^5$ has 5 unpaired electrons,
$\ce{Mn^3+}$ $[\mathrm{Ar}]\,3\mathrm d^4$ has 4: $\ce{Mn^2+}$ has more.
\end{solution}

\wq{4}{Beyond the known table}
The last element of period 7 is oganesson, Og ($Z = 118$). (a) Write the valence configuration of
Og. (b) Using Madelung's rule, which subshells would period 8 fill, and in which order? (Remember that
$\ell = 4$ is called g.) (c) Predict the atomic number of the noble gas at the end of period 8.
\begin{solution}
(a) $[\mathrm{Rn}]\,5\mathrm f^{14} 6\mathrm d^{10} 7\mathrm s^2 7\mathrm p^6$.\\
(b) After 7p ($n+\ell = 8$): 8s ($n+\ell = 8$, larger $n$ than 7p), then $n+\ell = 9$ in order of
increasing $n$: 5g, 6f, 7d, 8p.\\
(c) $2 + 18 + 14 + 10 + 6 = 50$ elements, so the next noble gas would have $Z = 118 + 50 = 168$.
(Real super-heavy atoms are expected to deviate from this simple rule, but it is the Madelung
prediction.)
\end{solution}
